% ---------------------------------------------------------------
%  PHAN II -- NOI DUNG
%  Muc II.1 -- Co so li luan va co so phap li
% ---------------------------------------------------------------
\section{NỘI DUNG SÁNG KIẾN}

\subsection{Cơ sở lí luận và cơ sở pháp lí}

\subsubsection{Cơ sở pháp lí}

Sáng kiến được xây dựng trên các căn cứ pháp lí sau:

\begin{itemize}[leftmargin=1.2cm, itemsep=3pt]
  \item \textbf{Thông tư số 32/2018/TT-BGDĐT} ngày 26/12/2018 ban hành Chương
        trình Giáo dục phổ thông. Chương trình môn Toán xác định mục tiêu hình
        thành và phát triển năm thành tố của năng lực toán học, trong đó có
        năng lực giải quyết vấn đề toán học và năng lực sử dụng công cụ,
        phương tiện học toán. Đây là căn cứ để thiết kế câu hỏi theo mức độ
        năng lực chứ không theo mức độ ghi nhớ.

  \item \textbf{Thông tư số 22/2021/TT-BGDĐT} ngày 20/7/2021 quy định về đánh
        giá học sinh trung học cơ sở và trung học phổ thông. Thông tư yêu cầu
        việc đánh giá phải \textit{``vì sự tiến bộ của học sinh''}, coi trọng
        động viên, khuyến khích sự cố gắng của người học, bảo đảm kịp thời,
        công bằng, khách quan. Hai yêu cầu \textit{kịp thời} và \textit{công
        bằng} chính là hai vấn đề mà sáng kiến này hướng tới giải quyết bằng
        công nghệ: chấm ngay khi học sinh nộp bài, và mỗi học sinh một đề
        riêng.

  \item \textbf{Quyết định số 764/QĐ-BGDĐT} ngày 08/3/2024 của Bộ Giáo dục và
        Đào tạo về việc phê duyệt cấu trúc định dạng đề thi tốt nghiệp trung
        học phổ thông từ năm 2025, với ba dạng câu hỏi trắc nghiệm: nhiều lựa
        chọn, đúng/sai và trả lời ngắn. Đây là căn cứ để hệ thống sinh đủ bốn
        dạng câu hỏi (ba dạng trắc nghiệm nêu trên và tự luận).

  \item \textbf{Quyết định số 131/QĐ-TTg} ngày 25/01/2022 của Thủ tướng Chính
        phủ phê duyệt Đề án ``Tăng cường ứng dụng công nghệ thông tin và
        chuyển đổi số trong giáo dục và đào tạo''.

  \item \textbf{Nghị quyết số 57-NQ/TW} ngày 22/12/2024 của Bộ Chính trị về
        đột phá phát triển khoa học, công nghệ, đổi mới sáng tạo và chuyển đổi
        số quốc gia, trong đó xác định phải đẩy mạnh ứng dụng trí tuệ nhân tạo
        và dữ liệu lớn trong các lĩnh vực, đưa giáo dục trở thành một trong
        những lĩnh vực được ưu tiên chuyển đổi số.

  \item \textbf{Nghị định số 13/2012/NĐ-CP} ngày 02/3/2012 của Chính phủ ban
        hành Điều lệ Sáng kiến và các văn bản hướng dẫn hiện hành của Sở Giáo
        dục và Đào tạo về công tác sáng kiến năm học \NAMHOC.
\end{itemize}

\ghichu{Thay dòng cuối bằng đúng số hiệu công văn hướng dẫn sáng kiến năm nay
của Sở mình --- hội đồng chấm thường đối chiếu ngay dòng này.}

\subsubsection{Cơ sở lí luận}

\textbf{a) Đánh giá vì sự tiến bộ của người học đòi hỏi phản hồi tức thời}

Lí luận dạy học hiện đại phân biệt \textit{đánh giá kết quả học tập}
(\textit{assessment of learning}) với \textit{đánh giá vì học tập}
(\textit{assessment for learning}). Ở dạng thứ hai, giá trị của bài kiểm tra
không nằm ở điểm số mà nằm ở thông tin phản hồi mà học sinh nhận được và ở
việc học sinh được luyện lại ngay phần mình còn yếu. Muốn vậy phải đáp ứng hai
điều kiện: phản hồi phải đến \textit{ngay} sau khi làm bài, và học sinh phải
có \textit{đủ đề} để luyện lại nhiều lần cùng một dạng toán. Cách làm thủ công
không đáp ứng được cả hai điều kiện này --- giáo viên không thể chấm ngay cho
mấy trăm bài, cũng không thể soạn cho mỗi em một đề luyện tập riêng.

\textbf{b) Câu hỏi tham số hoá và tính tương đương của đề}

Trong đo lường giáo dục, hai câu hỏi được coi là \textit{tương đương} khi
chúng cùng đo một chỉ báo năng lực, cùng dạng thao tác tư duy và cùng mức độ
nhận thức, dù số liệu cụ thể khác nhau. Từ nhận xét đó hình thành kĩ thuật
\textit{câu hỏi tham số hoá} (\textit{parameterized item} hay
\textit{algorithmic item}): thay vì lưu một câu hỏi cố định, người ta lưu quy
tắc sinh ra câu hỏi cùng với quy tắc sinh đáp án và lời giải. Mỗi lần gọi, quy
tắc cho ra một câu hỏi mới với bộ số liệu mới nhưng vẫn giữ nguyên dạng toán
và độ khó.

Kĩ thuật này giải quyết triệt để vấn đề công bằng trong kiểm tra: các mã đề
không còn là một đề hoán vị thứ tự mà là những đề thực sự khác nhau về số
liệu. Đồng thời nó biến kho câu hỏi từ tài nguyên \textit{hữu hạn} thành tài
nguyên \textit{không cạn}: từ một câu gốc có thể sinh ra hàng trăm câu luyện
tập cho học sinh.

Điều kiện để áp dụng được kĩ thuật này là bộ tham số phải được ràng buộc chặt
để mọi biến thể đều có nghiệm đẹp, đúng phạm vi chương trình và không rơi vào
trường hợp suy biến --- đây chính là phần công sức chuyên môn của giáo viên bộ
môn, máy không làm thay được.

\textbf{c) Ma trận đề và bốn mức độ nhận thức}

Đề kiểm tra chỉ có giá trị đo lường khi được xây dựng từ một ma trận xác định
rõ tỉ lệ câu hỏi theo từng đơn vị kiến thức và từng mức độ nhận thức: nhận
biết, thông hiểu, vận dụng, vận dụng cao. Về nguyên tắc, trọng số của một đơn
vị kiến thức trong đề phải tương ứng với thời lượng dạy học của đơn vị đó
trong phân phối chương trình. Nguyên tắc này thường bị bỏ qua khi soạn đề thủ
công vì phải tính toán thủ công quá nhiều; ngược lại, đó lại là việc máy tính
làm rất tốt và tuyệt đối chính xác.

\textbf{d) Trí tuệ nhân tạo là trợ lí, không phải người thay thế giáo viên}

Các mô hình ngôn ngữ lớn hiện nay rất mạnh ở việc hiểu ngôn ngữ tự nhiên,
nhưng bản chất là mô hình dự đoán chuỗi kí tự có xác suất cao nhất, không phải
hệ suy luận toán học. Khi được yêu cầu tự ra đề và giải, mô hình có thể cho
lời giải trôi chảy nhưng sai ở bước tính toán, và lần chạy sau lại cho kết quả
khác lần trước. Một đề kiểm tra lấy điểm chính thức không chấp nhận được rủi
ro đó.

Vì vậy, sáng kiến này xác định rõ ranh giới phân vai: trí tuệ nhân tạo chỉ làm
những việc \textit{cần hiểu ngôn ngữ tự nhiên hoặc cần văn phong tự nhiên};
mọi việc có quy tắc rõ ràng --- tra bảng phân phối chương trình, dựng ma trận,
chọn mã câu hỏi, sinh số liệu, tính điểm, so khớp đáp án --- đều do mã nguồn
thực hiện. Còn quyền quyết định cuối cùng về nội dung và mức độ đề vẫn thuộc
về giáo viên. Nguyên tắc này được ghi thành quy tắc phát triển bắt buộc của hệ
thống: \textit{``Ưu tiên mã nguồn hơn trí tuệ nhân tạo; không để nghiệp vụ nằm
trong câu lệnh gửi cho mô hình.''}

\begin{figure}[H]
\centering
\begin{tikzpicture}[
  font=\small,
  box/.style={draw, rounded corners=3pt, align=center, minimum height=1.1cm,
              text width=3.3cm, thick},
  ai/.style={box, fill=blue!8},
  code/.style={box, fill=green!8},
  gv/.style={box, fill=orange!12},
  mui/.style={-latex, thick}
]
\node[gv] (gv) at (0,0) {\textbf{Giáo viên}\\[2pt]\scriptsize quyết định nội dung, mức độ, duyệt đề};
\node[ai] (ai) at (5,0) {\textbf{Trí tuệ nhân tạo}\\[2pt]\scriptsize hiểu yêu cầu bằng tiếng Việt};
\node[code] (code) at (10,0) {\textbf{Mã nguồn}\\[2pt]\scriptsize chọn câu, sinh số liệu, tính điểm};
\draw[mui] (gv) -- (ai);
\draw[mui] (ai) -- (code);
\node[align=center, font=\scriptsize\itshape] at (5,-1.6)
  {chỉ khâu ngôn ngữ};
\node[align=center, font=\scriptsize\itshape] at (10,-1.6)
  {mọi khâu có quy tắc};
\end{tikzpicture}
\caption{Phân vai giữa giáo viên, trí tuệ nhân tạo và mã nguồn trong hệ thống}
\end{figure}

\textbf{e) Một hệ thống phục vụ đồng thời hai đối tượng}

Từ bốn cơ sở trên, hệ thống được thiết kế để phục vụ đồng thời giáo viên và
học sinh, với hai vai trò khác nhau nhưng dùng chung một ngân hàng câu hỏi:

\begin{itemize}[leftmargin=1.2cm, itemsep=3pt]
  \item \textbf{Về phía học sinh:} chủ động tạo đề luyện tập đúng phạm vi mình
        cần bằng một câu yêu cầu tiếng Việt, làm bài trực tuyến và nhận điểm
        cùng đáp án ngay. Đây là điều kiện để đánh giá thường xuyên thực sự
        diễn ra \textit{vì sự tiến bộ} của các em.
  \item \textbf{Về phía giáo viên:} có nguồn đề đúng ma trận để dùng cho kiểm
        tra trên lớp, giao đề và trả điểm qua Google Classroom, đồng thời theo
        dõi được điểm trung bình và những bài học sinh hay sai để điều chỉnh
        việc dạy.
\end{itemize}

Cách phân vai này cũng phù hợp với triết lí phát triển đã đặt ra ngay từ đầu
của hệ thống: \textit{``Không dùng trí tuệ nhân tạo để thay thế giáo viên.
Trí tuệ nhân tạo đóng vai trò trợ lí. Giáo viên vẫn là người quyết định.''}

\clearpage
